\section{Visão Geral da Estatística e Análise de Dados}

\subsection{Uma Breve História da Estatística}

A palavra estatística deriva dos termos neolatim \textit{statisticum collegium} e italiano
\textit{statista}, que significam Colégio de Estado e Estadista respectivamente. A palavra
alemã \textit{Statistik} foi usada modernamente por Gottfried Achenwall em meados do
século XVIII para designar a coleta de dados políticos, econômicos e sociais de um país,
cujo equivalente em inglês era \textit{political arithmetic}. Posteriormente, através de
Sir John Sinclair e seus trabalhos no \textit{Statistical Accounts of Scotland}, a palavra
ganhou uma conotação mais ampla sobre coleta e classificação de dados.

Como vistos pelas suas raízes etimológicas, a estatística surgiu enquanto a disciplina
que se preocupava com a coleta, tratamento, classificação e análise de dados sobre
características político-sócio-econômicas de um país a fim de apoiar a formulação de
políticas públicas.

A Estatística Matemática tem seus fundamentos no desenvolvimento da Teoria das
Probabilidades, antes, com tratamento mais informal, nos trabalhos de Pierre de Fermat,
Blaise Pascal e Christian Huygens no século XVII e Thomas Bayes no século XVIII.
O rigor matemático também se expandiu no século XVIII com as obras \textit{Ars Conjectandi}
(1713) de Jakob Bernoulli e \textit{The Doctrine of Chances} de Abraham Moivre (1718).
A formalização matemática rigorosa das probabilidades veio ao auge com Andrei Kolmogorov
no século XIX.

A disciplina estatística também tem laços com a Teoria de Erros, cujas origens remontam
a \textit{Opera Miscellanea} de Roger Cottes de 1772, e sua reedição por Thomas Simpson
em 1755, que aplicou a teoria a erros de observação. A Teoria de Erros ganhou um aspecto
mais analítico ainda no final do século XVIII e ao longo do XIX nas mãos de Pierre-Simon
de Laplace, Joseph Louis Lagrange e Daniel Bernoulli. O seminal Método dos Mínimos Quadrados
(MMQ) foi desenvolvido e publicado independentemente por Adrien-Marie Legendre, Carl
Friedrich Gauss e Robert Adrain.

Com o desenvolvimento e sofisticação das técnicas estatísticas, como a amostragem,
modelagem e inferência, a estatística na virada do século XIX para o XX, expandiu-se
da esfera do Estado para se tornar uma disciplina aplicada em ciências naturais,
sociais, administração pública e privada. Neste contexto, figuram nomes George Box,
Pafnuty Chebyshev, Gertrude Cox, George Dantzig, Edward Deming, Bruno de Finetti,
Sir Ronald Fisher, Sir Francis Galton, Aleksandr Lyapunov, Karl Pearson, Charles Spearman
e John Tukey.

\subsection{Modelos Estatísticos}

A estatística é um ramo da matemática aplicada e da metodologia científica
que busca extrair, reduzir, analisar e modelar um conjunto de dados que
amostram uma população. Para isso, são usadas principalmente as ferramentas
teóricas providas pela Teoria das Probabilidades, pois entende-se que os
dados originam-se de fenômenos aleatórios, e, uma vez modelados, geralmente
dois objetivos visam ser cumpridos: estimar ou predizer o comportamento dos
dados no futuro, ou realizar inferência, que consiste em detectar os padrões
embutidos. A inferência pode ser feita por duas abordagens: a inferência
dedutiva, que aceita determinadas premissas para chegar às conclusões, ou
a inferência indutiva, que parte de casos particulares para generalizar.

Ainda sobre a inferência estatística, quando tenta-se modelar um determinado
conjunto de dados, tentamos identificar padrões e tendências de comportamento
através de modelos estatísticos. Supondo que temos um conjunto de dados $D$,
e identificamos que conseguimos ajustar a eles um modelo $M$, teremos então que
\[
	D = M + R
\]
onde $R$ descreve a parte aleatória dos dados que não pode ser capturada pelo
modelo, chamada de resíduos. Normalmente, busca-se que resíduos não devam conter nenhuma
suavidade, pois isso implicaria em padrões que o modelo falhou em capturar, e mais
suavização por parte de $M$ será necessária.

\subsection{Análise de Dados}

A Análise de Dados consiste na inspeção, limpeza, transformação e modelagem de dados
com o objetivo claro de obter informações úteis, informar conclusões ou apoiar a
tomada de decisão. Em resumo, ela é o processo de extração de dados brutos para
geração de conhecimento acionável, geralmente feito num ciclo iterativo. As fases desse
processo podem ser divididas em:

\begin{enumerate}
	\item Requisitos de Dados: a escolha de dados a serem extraídos dependem das
	      análises que se deseja realizar. Aqui define-se a \textbf{unidade experimental},
	      que pode consistir de um indivíduo ou uma população deles, e as \textbf{variáveis}
	      específicas dessa unidade, como idade e renda para uma população de pessoas.

	\item Coleta de Dados: consiste na tarefa de obter os dados de diferentes fontes,
	      como bancos de dados estruturados ou não-estruturados, sensores físicos, formulários,
	      uma interface de programação de aplicação (API), arquivos binários ou de textos, entre
	      muitas outras fontes.

	\item Processamento e Limpeza: inclui tarefas de organizar os dados em formatos
	      estruturados, identificar tipos, deduplicar, tratar valores faltantes, incorretos
	      e desviantes e transformar o formato de dados para outros mais convenientes.

	\item Análise Exploratória: compreende uma gama de técnicas para descrever e resumir
	      os dados obtidos à procura de informações que ajudem entender relações entre as
	      variáveis e guie às técnicas a serem aplicadas nas fases posteriores. Aqui
	      é comumente aplicadas técnicas da Estatística Descritiva e métodos gráficos.

	\item Modelagem: são usados algoritmos e modelos estatísticos para extrair padrões
	      latentes nos dados, muitas vezes não triviais de serem identificados. As
	      técnicas de aprendizado estatístico são especialmente úteis nessa etapa.

	\item Comunicação e Produto de Dados: a depender do contexto e das finalidades definidas
	      a priori, busca-se maneiras de comunicar esses dados aos atores interessados na
	      forma de visualizações e relatórios. Nos últimos anos, a disciplina de
	      \textbf{Narrativa de Dados} tem se destacado para realizar essa comunicação.
\end{enumerate}

A Análise de Dados, do ponto de vista da Estatística, poderia ser dividida em
\begin{enumerate}
	\item Estatística Descritiva: utilizar medidas de resumo numéricas para os dados
	      a fim de descrever padrões visíveis, geralmente respondendo perguntas do tipo
	      ``o que aconteceu?''.

	\item Análise Exploratória: semelhante a anterior, porém empregando métodos mais
	      flexíveis e visuais que permitam identificar padrões mais ocultos, correlações e
	      anomalias. Ocorre num processo iterativo de formulação de perguntas, exploração
	      e descobertas.

	\item Análise Confirmatória: aplica as técnicas da inferência estatística para
	      confirmar ou rejeitar hipóteses existentes das etapas anteriores
\end{enumerate}
